\section{Numerical experiments}
\label{sec:numerical-experiments}

We conduct a compact simulation to assess the finite-sample behavior of the
six independence tests.  The experiment is deliberately targeted at two
questions raised by the preceding theory: whether effective-sample-size
calibration remains adequate under moderately unequal exogenous weights, and
whether the general-dependence statistics detect a symmetric nonmonotone
alternative that is difficult for signed or quadrant-based coefficients.  The
complete design is implemented in
\texttt{paper/simulation/run-study.R} using version 0.3.0 of the R package.

\subsection{Design}

For each design point, predictors and innovations are independent standard
normal random variables.  Under the null, $Y=\varepsilon$.  The two
alternatives are
\begin{align*}
  Y &= 0.25X + \sqrt{1-0.25^2}\,\varepsilon,
      &&\text{(linear)}, \\
  Y &= \frac{0.5(X^2-1)}{\sqrt{2}}
       + \sqrt{1-0.5^2}\,\varepsilon,
      &&\text{(quadratic)}.
\end{align*}
The linear design has correlation $0.25$.  In the quadratic design the
standardized quadratic component and the response both have unit variance,
while all signed first-order association vanishes by symmetry.

We compare equal weights with the deterministic profile
\begin{align*}
  w_i
  &= \exp\left\{0.8\Phi^{-1}
       \left(\frac{i-1/2}{n}\right)\right\},
  & i&=1,\ldots,n.
\end{align*}
The profile approximates log-normal weights but is fixed independently of the
observations, exactly matching the exogenous-weight setting of
Section~\ref{sec:independence-tests}.  We consider nominal sample sizes
$n\in\{50,100,200,500\}$.  Null rejection probabilities use 20000 replications
at each combination of sample size and weight profile; power uses 10000
replications.  Pearson, Spearman, Kendall, Blomqvist, and Hoeffding tests are
two-sided, while the Chatterjee test uses its natural upper-tailed alternative.
All variables are continuous because the current Hoeffding implementation does
not support ties.  The complete study ran single-threaded in 452 seconds on
the reference machine.

\subsection{Null calibration}

\begin{table}[t]
  \centering
  \small
  \setlength{\tabcolsep}{3.5pt}
  \caption{Empirical rejection probabilities under independence at nominal level $0.05$. The Monte Carlo standard error is at most $0.0049$.}
  \label{tab:simulation-size}
  \begin{tabular}{rrcccccc}
    \toprule
    $n$ & $n_{\mathrm{eff}}$ & Pear. & Spear. & Kend. & Blom. & Hoef. & Chatt. \\
    \midrule
    \multicolumn{8}{l}{Equal weights} \\
    50 & 50.0 & 0.046 & 0.041 & 0.047 & 0.096 & 0.051 & 0.056 \\
    100 & 100.0 & 0.050 & 0.042 & 0.051 & 0.071 & 0.051 & 0.055 \\
    200 & 200.0 & 0.049 & 0.040 & 0.047 & 0.063 & 0.046 & 0.048 \\
    500 & 500.0 & 0.056 & 0.046 & 0.052 & 0.054 & 0.053 & 0.052 \\
    \addlinespace
    \multicolumn{8}{l}{Unequal weights} \\
    50 & 28.5 & 0.037 & 0.037 & 0.040 & 0.072 & 0.067 & 0.041 \\
    100 & 55.6 & 0.047 & 0.044 & 0.052 & 0.067 & 0.072 & 0.057 \\
    200 & 109.3 & 0.040 & 0.040 & 0.044 & 0.056 & 0.052 & 0.040 \\
    500 & 269.1 & 0.049 & 0.035 & 0.041 & 0.043 & 0.046 & 0.051 \\
    \bottomrule
  \end{tabular}
\end{table}


Table~\ref{tab:simulation-size} shows that most rejection probabilities are
close to the nominal level.  The entries are to be read against a maximum
Monte Carlo standard error of $0.0020$.  Unequal weighting does not create a
systematic size distortion, although its nominal sample sizes correspond to
substantially smaller effective samples.  Blomqvist's asymptotic test is
liberal in the smallest equal-weight samples, with rejection probability
$0.092$ at $n=50$; it improves as $n$ grows, but remains mildly liberal at
$0.061$ for $n=500$.  Unequal-weight Hoeffding is mildly liberal at $n=50$ and
$100$, with rejection probabilities $0.066$ and $0.060$, and is close to
nominal from $n=200$ onward.  Spearman's test is mildly conservative
throughout, with rejection probabilities between $0.037$ and $0.044$, which is
consistent with the factor $1.06$ discussed in
Section~\ref{subsec:standard-tests}.  Apart from Blomqvist's test, all entries
at $n=500$ lie between $0.042$ and $0.053$.

\subsection{Power}

\begin{table}[p]
  \centering
  \small
  \setlength{\tabcolsep}{3.5pt}
  \caption{Empirical power at nominal level $0.05$. The Monte Carlo standard error is at most $0.0158$.}
  \label{tab:simulation-power}
  \begin{tabular}{rrcccccc}
    \toprule
    $n$ & $n_{\mathrm{eff}}$ & Pear. & Spear. & Kend. & Blom. & Hoef. & Chatt. \\
    \midrule
    \multicolumn{8}{l}{Linear alternative, equal weights} \\
    50 & 50.0 & 0.427 & 0.363 & 0.374 & 0.303 & 0.354 & 0.102 \\
    100 & 100.0 & 0.697 & 0.623 & 0.644 & 0.410 & 0.585 & 0.147 \\
    200 & 200.0 & 0.957 & 0.930 & 0.934 & 0.662 & 0.893 & 0.189 \\
    500 & 500.0 & 1.000 & 1.000 & 1.000 & 0.963 & 0.999 & 0.335 \\
    \addlinespace
    \multicolumn{8}{l}{Linear alternative, unequal weights} \\
    50 & 28.5 & 0.243 & 0.216 & 0.223 & 0.157 & 0.252 & 0.089 \\
    100 & 55.6 & 0.470 & 0.405 & 0.418 & 0.238 & 0.395 & 0.123 \\
    200 & 109.3 & 0.753 & 0.677 & 0.707 & 0.391 & 0.654 & 0.139 \\
    500 & 269.1 & 0.988 & 0.977 & 0.981 & 0.732 & 0.956 & 0.230 \\
    \addlinespace
    \multicolumn{8}{l}{Quadratic alternative, equal weights} \\
    50 & 50.0 & 0.134 & 0.065 & 0.077 & 0.091 & 0.216 & 0.310 \\
    100 & 100.0 & 0.144 & 0.051 & 0.063 & 0.071 & 0.469 & 0.486 \\
    200 & 200.0 & 0.179 & 0.060 & 0.080 & 0.067 & 0.927 & 0.779 \\
    500 & 500.0 & 0.174 & 0.066 & 0.086 & 0.065 & 1.000 & 0.975 \\
    \addlinespace
    \multicolumn{8}{l}{Quadratic alternative, unequal weights} \\
    50 & 28.5 & 0.123 & 0.059 & 0.074 & 0.084 & 0.150 & 0.193 \\
    100 & 55.6 & 0.142 & 0.068 & 0.085 & 0.066 & 0.282 & 0.336 \\
    200 & 109.3 & 0.145 & 0.070 & 0.091 & 0.084 & 0.556 & 0.535 \\
    500 & 269.1 & 0.129 & 0.055 & 0.073 & 0.066 & 0.987 & 0.843 \\
    \bottomrule
  \end{tabular}
\end{table}


The linear design in Table~\ref{tab:simulation-power} favors the classical
correlation and monotone rank tests.  With equal weights, Pearson power grows
from $0.433$ at $n=50$ to $0.950$ at $n=200$, followed closely by Spearman and
Kendall.  Chatterjee's coefficient is less sensitive here because its
population signal is of second order in a weak linear association.  The
ordering changes for the quadratic alternative: at $n=200$ with equal weights,
Hoeffding's $D$ and Chatterjee's $\xi$ attain power $0.917$ and $0.755$, while
Spearman, Kendall, and Blomqvist remain between $0.065$ and $0.082$ because
their population coefficients vanish by symmetry.
Pearson's test also has a nonzero rejection probability despite zero
covariance because its usual Gaussian calibration is not designed for this
form of nonlinear dependence.  As expected, weight concentration reduces
power in both designs by lowering the effective sample size.  Thus the
experiment supports the main qualitative claims without suggesting that any
single coefficient is uniformly preferable.
